\section{A flux theorem and a Lean-verified six-boundary obstruction}
For any vertex set $W$ and target $t$, internal edges cancel to give
\begin{equation}\label{eq:flux}
 \sum_{x\in W}\stab(x,t)=
 \sum_{xy\in E:\,x\in W,\,y\notin W}(d(x,t)-d(y,t)).
\end{equation}
Strictly positive outgoing flux therefore certifies a boundary vertex in $W$.
Say a root $u$ has the \emph{unique first-step property} if every $x\ne u$
has exactly one neighbor $v_i$ of $u$ with $d(v_i,x)=d(u,x)-1$.
Put $f(x)=d(u,x)$ and let $S_i$ be the sector of vertices with that first
neighbor. The sectors partition $V\setminus\{u\}$.

\begin{theorem}[Sector lower bound]\label{thm:sectors}
Let $u$ have the unique first-step property, and let $J$ be the set of
indices $i$ for which there is a nonadjacent root neighbor $v_j$ and
an $x\in S_i$ with $d(v_j,x)=f(x)$. Then
\[
 b(G)\ge\dd(u)+|J|+\ind_{\{u\in\bd G\}}.
\]
If $u$ is not a cut vertex and $G[N(u)]$ has $c\ge2$ components, then
\[
 b(G)\ge\dd(u)+c+\ind_{\{u\in\bd G\}}.
\]
In particular, if $u$ is not a cut vertex, independent neighbors and $\dd(u)\ge2$ give
$b(G)\ge2\dd(u)+\ind_{\{u\in\bd G\}}$.
\end{theorem}
\begin{proof}
An edge between distinct sectors preserves $f$: a rise of one would
produce a shortest path with a second first neighbor. Consequently the
outgoing $f$-flux of $S_i$ is exactly one, from the edge $v_i u$.
Each sector contains a boundary vertex.

For $i\in J$ choose $j$ as in the statement and partition
\[
 P=\{x\in S_i:d(v_j,x)=f(x)\},\qquad T=S_i\setminus P.
\]
The set $P$ is nonempty and $v_i\in T$, because $v_i,v_j$ are nonadjacent.
For an edge $xy$ from $P$ to $T$, distance Lipschitzness gives
$f(x)\ge f(y)$, since $d(v_j,y)=f(y)+1$.
Edges from $P$ to other sectors preserve $f$, and none goes to $u$.
A vertex of $P$ minimizing $f$ has a geodesic predecessor outside $P$,
giving a strictly positive contribution. Thus $P$ has positive $f$-flux.

For $T$ use target $v_j$. The edge $v_i u$ contributes $2-1=1$.
An edge $xy$ from $T$ to $P$ contributes $f(x)+1-f(y)\ge0$.
An edge leaving $S_i$ preserves $f$ and cannot enter $S_j$, since its
distances to $v_j$ would differ by two. In every other sector its endpoint
has distance at most $f(x)+1$ from $v_j$, giving a nonnegative contribution.
Thus $T$ also has positive flux. The two disjoint sets give two distinct
boundary vertices in every sector indexed by $J$, proving the first bound.

Let $Q=G[N(u)]$, and let $H$ join two root neighbors when their sectors
have a crossing edge. Then $Q\subseteq H$, and each endpoint of an edge
of $H\setminus Q$ belongs to $J$. If $G-u$ is connected, projecting a
path onto its sectors shows that $H$ is connected. Each of the $c$
components of $Q$ is incident with an edge leaving it, so $|J|\ge c$.
\end{proof}

\begin{theorem}[Degree-three obstruction]\label{thm:degree3}
If a graph has a non-cut interior vertex of degree three, it has at least
six boundary vertices. The constant six is sharp.
\end{theorem}
\begin{proof}
At an interior degree-three root, every nonroot target has at least one
closer neighbor. Two closer neighbors would give positive stability,
since at most one could be farther. Thus there is exactly one closer
neighbor, and at least one farther neighbor: the unique first-step property
holds.
Its three neighbors are independent. A path of length two among them
would give its middle neighbor a positive score at the root. If there is
exactly one neighbor edge $v_1v_2$, then for every $x\in S_1$ the
distance from $v_2$ equals $f(x)$. Interiorness forces the distance from
$v_3$ to equal $f(x)+1$. No edge can therefore join $S_1$ to $S_3$;
the same holds for $S_2,S_3$. This contradicts connectivity of $G-u$.
Theorem~\ref{thm:sectors} now gives six.

For sharpness, take three internally disjoint paths between two terminals
of lengths $1,4,4$. The six internal degree-two vertices are boundary;
both degree-three terminals have stability zero. Hence $b=6$.
\end{proof}
The full statement of Theorem~\ref{thm:degree3}, including extraction of
six distinct actual boundary vertices from graph distances and connectivity
after deleting the root, is proved in the Lean file\\
\texttt{RootSectors.lean}.

